\section{Conclusiones y recomendación}
\begin{enumerate}
\item La mejora depende de organizar fuentes, responsabilidades y vigencia. Los terminales amplían el acceso presencial cuando se integran con ese proceso y mantienen una vía móvil de consulta. La finalidad se apoya en obligaciones de la Ley Universitaria y en estándares del modelo de acreditación SINEACE 2025, además de la necesidad declarada.
\item El aporte TI comprende publicación trazable, distribución, interacción, control de acceso, continuidad y recuperación. Cada tipo de aviso tiene una fuente oficial y un emisor propuesto. La aceptación se basará en pruebas del servicio integrado y en el desempeño observado con usuarios.
\item El requerimiento se expresa en diez especificaciones sin marca ni modelo, como exige el Reglamento de la Ley N.° 32069. Cada cantidad y capacidad se justifica frente a valores menores y mayores. Las marcas sólo sirven para valorizar, mediante cuadros comparativos con proveedores identificados y fichas registradas en el anexo C. Se excluyen Windows 10, los procesadores no admitidos por Windows 11 y el cable de aluminio recubierto.
\item El subsistema energético combina 3,69 kWp, un inversor con 16 A por seguidor y operación hasta 4\,000 m, y una batería de 4,8 kWh que entrega 2,76 kWh de diseño. El precio más bajo no definió el inversor: tres de las opciones comparadas, dos de ellas más baratas, limitan la corriente de los módulos, y una tampoco opera a la altitud de Cusco. Protecciones, conductores, tierra y caídas de tensión tienen cálculo explícito. La red de siete enlaces y la videovigilancia de 30 días protegen los equipos con el menor dimensionamiento que cumple.
\item El escenario integral A2 requiere S/ 93\,417,69 de inversión con reservas, con 60\% de su base respaldada por precios publicados. A1 mantiene el servicio y respaldo comparables por S/ 83\,287,25. La decisión solar debe justificar el incremento de S/ 10\,130,44: el ahorro eléctrico del horizonte no basta por sí solo y, en valor actual al 8\%, A2 supera a A1 en S/ 7\,245,86. Por estudiante de la población objetivo, A1 cuesta S/ 107,49 más que A0 en cinco años y A2 añade S/ 8,52 sobre A1.
\item La ejecución, transferencia, mantenimiento por equipo y evaluación se organizan mediante entregables, responsables, recursos, pruebas y puntos de decisión. El seguimiento de operación determinará la mejora efectiva y la conveniencia de expansión.
\end{enumerate}

Se recomienda confirmar la línea base y la matrícula oficial, y demostrar el flujo editorial y el cliente antes de contratar. Con esas evidencias se cerrará la comparación A0/A1/A2. La ampliación por edificios requerirá capacidad institucional, costo de ciclo de vida y resultados del piloto, conservando la misma disciplina de especificación, selección y aceptación.
